\chapter{项目一：命令行词频统计器}

\section{项目目标}

这一章把前面学过的内容放进一个小工具：读取文本文件，统计单词出现次数，输出出现次数最高的前 \code{k} 个单词。这个项目覆盖文件 RAII、字符串处理、哈希表、排序、错误处理和命令行参数。

目标命令如下：

\begin{lstlisting}[language=bash,caption={词频工具目标用法}]
wordfreq --input article.txt --top 20
\end{lstlisting}

输出示例：

\begin{lstlisting}[caption={词频输出示例}]
the 128
system 42
memory 31
\end{lstlisting}

\section{第一版：把文件读进字符串}

最简单的做法是逐词读取：

\begin{lstlisting}[language=C++,caption={逐词读取文件}]
#include <fstream>
#include <string>
#include <vector>

std::vector<std::string> read_words(const std::string& path)
{
    std::ifstream input{path};
    std::vector<std::string> words;
    std::string word;
    while (input >> word) {
        words.push_back(word);
    }
    return words;
}
\end{lstlisting}

这段代码利用 \code{std::ifstream} 的 RAII 自动关闭文件。但它还不够：文件打不开时返回空数组，调用者无法区分“空文件”和“文件不存在”；标点也没有处理，\code{"system"} 和 \code{"system,"} 会被当成两个词。

\section{定义错误和清洗规则}

先定义读取结果：

\begin{lstlisting}[language=C++,caption={读取结果类型}]
struct WordReadError {
    std::string message;
};

struct WordReadResult {
    std::vector<std::string> words;
    std::optional<WordReadError> error;
};
\end{lstlisting}

单词清洗规则先保持简单：保留字母和数字，统一转小写，其他字符当作分隔。

\begin{lstlisting}[language=C++,caption={清洗一行文本}]
#include <cctype>

std::vector<std::string> split_words(std::string_view line)
{
    std::vector<std::string> words;
    std::string current;
    for (unsigned char ch : line) {
        if (std::isalnum(ch)) {
            current.push_back(static_cast<char>(std::tolower(ch)));
            continue;
        }
        if (!current.empty()) {
            words.push_back(current);
            current.clear();
        }
    }
    if (!current.empty()) {
        words.push_back(current);
    }
    return words;
}
\end{lstlisting}

这里用 \code{unsigned char} 是细节。很多 \code{ctype} 函数要求参数能表示为 \code{unsigned char} 或 EOF，直接传负的 \code{char} 可能出问题。

\section{统计和排序}

统计用 \code{unordered_map}，排序用 \code{vector}。

\begin{lstlisting}[language=C++,caption={统计词频}]
#include <unordered_map>

using WordCounts = std::unordered_map<std::string, int>;

WordCounts count_words(std::span<const std::string> words)
{
    WordCounts counts;
    for (const std::string& word : words) {
        ++counts[word];
    }
    return counts;
}
\end{lstlisting}

为了排序，把哈希表转成数组：

\begin{lstlisting}[language=C++,caption={取 top k 高频词}]
#include <algorithm>

struct WordCount {
    std::string word;
    int count = 0;
};

std::vector<WordCount> top_words(const WordCounts& counts, std::size_t limit)
{
    std::vector<WordCount> items;
    items.reserve(counts.size());
    for (const auto& [word, count] : counts) {
        items.push_back(WordCount{.word = word, .count = count});
    }

    std::ranges::sort(items, [](const WordCount& lhs, const WordCount& rhs) {
        if (lhs.count != rhs.count) {
            return lhs.count > rhs.count;
        }
        return lhs.word < rhs.word;
    });

    if (items.size() > limit) {
        items.resize(limit);
    }
    return items;
}
\end{lstlisting}

排序规则必须完整：次数高的排前面；次数相同按字典序。这让输出稳定，测试也稳定。

\section{复杂度和替代方案}

设单词总数为 \code{n}，不同单词数为 \code{m}。统计平均复杂度是 \complexity{O(n)}，排序是 \complexity{O(m log m)}。如果 \code{k} 很小而 \code{m} 很大，可以用小顶堆把排序降到 \complexity{O(m log k)}。但第一版项目更重视正确性和可读性，全排序完全可以接受。

\section{测试清单}

\begin{itemize}
  \item 空文件：输出为空。
  \item 大小写：\code{Cpp} 和 \code{cpp} 合并。
  \item 标点：\code{memory, memory.} 计为同一个词。
  \item 次数相同：按字典序输出。
  \item 文件不存在：返回明确错误。
\end{itemize}

\begin{keyidea}
一个小工具也要有结构：输入和清洗、统计、排序、输出分开。这样每一步都能单独测试，后续替换算法也不会牵动整段程序。
\end{keyidea}

\section{把需求拆成可测试流水线}

词频工具看起来是一个命令，其实可以拆成五步：

\begin{enumerate}
  \item 解析命令行，得到输入路径和 \code{top} 数量。
  \item 读取文件，每次拿到一行文本。
  \item 把一行文本清洗成单词。
  \item 统计所有单词出现次数。
  \item 排序、截断并输出。
\end{enumerate}

拆分的标准是：每一步都有清楚输入输出，可以单独测试。比如 \code{split_words("Cpp, CPP!")} 应该得到两个 \code{cpp}；这个测试不需要真实文件。等清洗规则可靠后，再测试文件读取。

\section{读取文件时保留错误}

前面的 \code{read_words} 会把文件不存在和空文件都变成空数组。改进版应该显式检查文件打开：

\begin{lstlisting}[language=C++,caption={逐行读取并返回错误}]
WordReadResult read_words_from_file(const std::string& path)
{
    std::ifstream input{path};
    if (!input) {
        return {.error = WordReadError{"cannot open input file"}};
    }

    WordReadResult result;
    std::string line;
    while (std::getline(input, line)) {
        std::vector<std::string> words = split_words(line);
        result.words.insert(result.words.end(),
                            std::make_move_iterator(words.begin()),
                            std::make_move_iterator(words.end()));
    }
    if (input.bad()) {
        return {.error = WordReadError{"failed while reading input file"}};
    }
    return result;
}
\end{lstlisting}

这里用 \code{std::getline} 是为了按行清洗标点；用 \code{make_move_iterator} 是因为临时 \code{words} 里的字符串已经不再需要，可以移动到结果里。即便如此，第一版仍然把所有单词存进内存。若文件很大，下一步可以改成边读边统计。

\section{流式版本减少内存}

更好的结构是读取一行就更新计数，不必保存所有单词：

\begin{lstlisting}[language=C++,caption={边读边统计}]
struct CountResult {
    WordCounts counts;
    std::optional<WordReadError> error;
};

CountResult count_words_in_file(const std::string& path)
{
    std::ifstream input{path};
    if (!input) {
        return {.error = WordReadError{"cannot open input file"}};
    }

    CountResult result;
    std::string line;
    while (std::getline(input, line)) {
        for (std::string& word : split_words(line)) {
            ++result.counts[word];
        }
    }
    if (input.bad()) {
        return {.error = WordReadError{"failed while reading input file"}};
    }
    return result;
}
\end{lstlisting}

这个版本的峰值内存主要由不同单词数决定，而不是所有单词总数。我们不是为了炫耀优化，而是从数据规模推出结构：如果文本很大，保存全部中间单词没有必要。

\section{主程序如何连接模块}

入口函数把前面步骤串起来：

\begin{lstlisting}[language=C++,caption={词频工具主流程}]
int run_wordfreq(const Options& options)
{
    const CountResult counted = count_words_in_file(options.input_path);
    if (counted.error.has_value()) {
        std::cerr << counted.error->message << '\n';
        return 2;
    }

    const std::vector<WordCount> top = top_words(counted.counts, options.top);
    for (const WordCount& item : top) {
        std::cout << item.word << " " << item.count << '\n';
    }
    return 0;
}
\end{lstlisting}

这个函数仍然很短，因为细节已经在可测试函数里。入口处理错误消息和输出，核心函数处理数据规则。以后如果要输出 JSON，只需要替换输出层；如果要支持停用词，只需要在清洗或统计之间加一步过滤。

\section{项目验收标准}

一个可交付的小工具不只是“我本机能跑”。至少要满足：

\begin{itemize}
  \item 命令行缺少 \code{--input} 或 \code{--top} 时给出错误。
  \item 文件打不开时返回非零退出码。
  \item 标点、大小写、空行处理符合文档规则。
  \item 次数相同时输出顺序稳定。
  \item 核心函数有测试，测试不依赖真实工作目录。
\end{itemize}

\begin{exercisebox}
给词频工具增加 \code{--min-length N} 参数，只统计长度至少为 \code{N} 的单词。先写清楚 \code{N = 0}、缺值、非数字、负数分别怎么处理，再改参数解析和统计流程。
\end{exercisebox}
